\begin{code}[hide]
module vqupit.sections.framework where

open import Data.Nat using (ℕ)
open import Data.Product using (_×_ ; _,_)
open import Level using (Level ; _⊔_ ; 0ℓ)
open import Function.Definitions using (Congruent ; Injective)
open import Algebra.Bundles using (Group)
open import Word.Base
open import Presentation.Definitions using (_IsPresentationOf_)
open import Presentation.Construct.Base using (_⋄_⋄_ ; ConjRelʷ ; EmptyRel)

infix 4 _≈_ _≈ₙ_ _≈₁_ _≈₂_ _~_ _===₁_ _===₂_

postulate
  proof-elided : ∀ {ℓ} {A : Set ℓ} → A
  X Y C Sem N H : Set
  c d : Level
  |NF| : Set
  ⟦_⟧ : Word X → Sem
  _≈_ : Word X → Word X → Set
  _≈ₙ_ : |NF| → |NF| → Set
  _≈₂_ : Sem → Sem → Set
  _≈₁_ : Word N → Word N → Set
  I : C
  f : X → Word Y
  h : C → Y → Word X × C
  [_] : C → Word Y
  _~_ : Word X × C → Word X × C → Set
  _===₁_ : WRel X
  _===₂_ : WRel Y
  Γ Δ : Set
  conj : H → N → Word N
  G1⋊G2 : Group 0ℓ 0ℓ
  RelTwist : Set → Set → Set
  _∪_ : Set → Set → Set

record NormalForm : Set where
  field inv-nf : |NF| → Word X
\end{code}
\section{The Framework, in Brief}\label{sec:framework}

The formalisation is built on an Agda library for presentations of
monoids and groups by generators and relations, normal forms, and
completeness of finite relation sets, with wire-indexed circuits as its
motivating instance~\cite{bian2026framework,qupit-code}.  We recall only
what the rest of the paper uses, with the design decisions that
matter: syntax is unquotiented, so normal forms and coset tables are
programs the evaluator can run; specifications are standard-library
function bundles; certificates are finite; and every construction comes
with a presentation theorem against a concrete group.

\subsection{Presented monoids as inductive congruences}\label{sec:presentations}

A word over a generator set $X$ is an unquotiented syntax tree with
constructors $[\_]\aid{ʷ}$, \aid{ε} and \aid{•}, deliberately not a
list, so that a word remembers its bracketing and rules apply to any
parenthesised subterm.  A relation set is $\aid{WRel}\;X =
\aid{Rel}\;(\aid{Word}\;X)\;0\ell$.  Given $\Gamma : \aid{WRel}\;X$,
\aid{{\_}==={\_}} always means the raw axioms and \aid{{\_}≈{\_}} the
monoid congruence they generate, an inductive family closed under
reflexivity, symmetry, transitivity, congruence of \aid{•},
associativity, the unit laws, and \aid{axiom}.  Thus $\approx$
\emph{is} the word problem of $\langle X \mid \Gamma\rangle$, and the
presented monoid is a setoid on words, not a quotient type: normal-form
maps, coset tables and sections are then structurally recursive
programs, and the sections into the syntax that the whole design rests
on have concrete words to land on.  Words have no primitive inverses, so
every presentation is at bottom a monoid presentation; a
\aid{Grouplike} witness (every generator has a left inverse up to
$\approx$) upgrades the presented monoid to the presented group.  The
extension of a generator map $f : X \to \aid{Word}\;Y$ to words is
written postfix, $(f\,\aid{ʷ})$, and a coset action $h : C \to Y \to
\aid{Word}\;X \times C$ extends to words as $(h\,\aid{ᵗ})$, threading
the coset from left to right.  A presentation of a concrete group is a
record, $\Gamma\;\aid{IsPresentationOf}\;G$, packaging a grouplike
witness with a semantics $\sem{\_} : \aid{Word}\;X \to G$ and a proof
that it is a group isomorphism from the presented group onto $G$.  Its
two halves are \emph{soundness}, that $\sem{\_}$ respects the rules,
and \emph{completeness}, that it is injective; every presentation
theorem in this paper proves soundness, completeness and surjectivity
separately, and the record assembles the isomorphism.

\subsection{Normal forms, and completeness by normalisation}\label{sec:nfproperty}

Normal-form witnesses are, on the nose, the standard library's function
bundles out of the word setoid: $\aid{NormalForm}\;\Gamma\;\NF$ is a
\aid{RightInverse}, the map \aid{nf} together with a \emph{section}
$\aid{inv-nf} : \NF \to \aid{Word}\;X$ and a proof $\aid{inv-nf} \circ
\aid{nf} \approx \mathrm{id}$.  The section is data: it is the algorithm
that realises each normal form as a canonical word.  Fix a semantics
$\sem{\_}$ into a setoid.  A normal form is \emph{unique for}
$\sem{\_}$ when representatives with equal denotations are equal, and
soundness plus a unique normal form yields completeness:

\begin{code}
record UniqueNormalForm (inv-nf : |NF| → Word X) : Set (c ⊔ d) where
  field
    unique : ∀ {u v : |NF|} → ⟦ inv-nf u ⟧ ≈₂ ⟦ inv-nf v ⟧ → u ≈ₙ v

by-normalization : {normalForm : NormalForm} → let open NormalForm normalForm in
  UniqueNormalForm inv-nf → Congruent _≈_ _≈₂_ ⟦_⟧ → Injective _≈_ _≈₂_ ⟦_⟧
\end{code}
\begin{code}[hide]
by-normalization = proof-elided
\end{code}

\noindent Given $\sem{w} \approx_2 \sem{v}$, soundness transports along
$w \approx \aid{inv-nf}\,(\aid{nf}\,w)$, uniqueness identifies the
normal forms, and injectivity of \aid{nf} closes the loop.  The value of
this five-line lemma is that every case study only ever proves
\aid{unique}, a statement about normal forms, i.e.\ about concrete
finite data, never about arbitrary words.  It is precisely the division
of labour of the qupit paper's Proposition~3.8 (uniqueness of the
symplectic normal form) and Proposition~4.3 (every circuit rewrites to
normal form), except that the rewriting half is a \emph{function} with
a verified section.  A converse, \aid{by-completeness}, recovers
\aid{UniqueNormalForm} from completeness plus an exact section; it is
how uniqueness witnesses are lifted through products.

\subsection{Circuits and the structural rules}\label{sec:circuits}

The circuit layer is parameterised by a gate family $\aid{Gate} :
\mathbb{N} \to \mathsf{Set}$ and defines wire-indexed generators
$\aid{Gen}\;n$ --- a $k$-ary gate on the bottom $k$ wires, at any width
$k + n$, or a generator shifted up one wire by \aid{{\_}↥} --- and
circuits $\aid{Circuit}\;n = \aid{Word}\;(\aid{Gen}\;n)$;
$c{\uparrow} = \aid{wmap}\;\aid{{\_}↥}$ shifts a circuit.  The indices are built from successors only, so that the coverage
checker solves its unification problems by injectivity of \aid{₁₊}
rather than by arithmetic, which is why coset tables can be defined by
direct pattern matching.  Given any gate-specific family of relations,
\aid{Lift-Relation} adds the structural rules that every circuit theory
shares --- the fixed-width shadow of the monoidal structure:
\aid{cong↑} (an equation may be shifted up), \aid{comm₁} and
\aid{comm₂} (gates on disjoint wires commute), and \aid{ω↑=ω} (a
scalar is its own shift).  The lemma \aid{lemma-cong↑} lifts the entire
congruence one wire up, so equational developments move freely along
the wire chain $\mathsf{Cir}\,k \le \mathsf{Cir}\,(k{+}1)$.

\subsection{Coset tables and towers}\label{sec:engine}

The library's normal forms follow one recipe from computational group
theory~\cite{sims1994computation,holt2005handbook}.  Given a chain $1 =
G_0 \le G_1 \le \dots \le G_n = G$ with a transversal $T_k$ of each
step, every element factors uniquely as $t_1 t_2 \cdots t_n$, a
mixed-radix numeral with one digit per level.  For circuits the chain is
handed to us by the wire structure: $\mathsf{Cir}\,k$ embeds into
$\mathsf{Cir}\,(k{+}1)$ as the circuits that do not touch the bottom
wire.  One level is a Reidemeister--Schreier situation: a subgroup
presentation $\Gamma$ and a group presentation $\Delta$, a finite type
$C$ of \emph{coset labels} with an identity coset $I$, a section $[\_]
: C \to \aid{Word}\;Y$ mapping each label to its representative, the
generator embedding $f$, and the \emph{coset table} $h : C \to Y \to
\aid{Word}\;X \times C$, which pushes one letter of $G$ past a coset
representative and returns the residual word of the subgroup and the
updated coset.  The classical algorithms \emph{find} such a table; the
library takes a claimed table and \emph{verifies} it against five
hypotheses:

\begingroup\renewcommand{\AgdaCodeStyle}{\scriptsize}
\begin{code}[hide]
module Transfer
\end{code}
\begin{code}
  (h=⁻¹f-gen : ∀ (x : X) → ([ x ]ʷ , I) ~ ((h ᵗ) I (f x)))                          -- 1
  (h-wd-ax : ∀ (c : C){u t : Word Y} → u ===₂ t → ((h ᵗ) c u) ~ ((h ᵗ) c t))        -- 2
  (f-wd-ax : ∀ {w v} → w ===₁ v → (f ʷ) w ≈₂ (f ʷ) v)                               -- 3
  ([I]≈ε : [ I ] ≈₂ ε)                                                              -- 4
  (h=ract : ∀ c b → let (b' , c') = h c b in [ c ] • [ b ]ʷ ≈₂ (f ʷ) b' • [ c' ])   -- 5
\end{code}
\endgroup
\begin{code}[hide]
  where
\end{code}

\noindent (1) the table inverts the embedding at the identity coset;
(2) the table respects $\Delta$'s axioms, coset by coset; (3) the
embedding respects $\Gamma$'s axioms; (4) the identity coset is
represented by $\varepsilon$; (5) the \emph{soundness law}: sliding one
letter past a representative matches the table.  From these the library
derives the normal-form transport $\aid{NormalForm}\;\Gamma\;\NF \to
\aid{NormalForm}\;\Delta\;(\NF \times C)$, and \aid{CosetTower}
iterates it up an $\mathbb{N}$-indexed family of presentations, one
level per wire.  Three of the hypotheses range over finite types and
two over axiom instances, so all five are case splits; in the warm-ups
of \S\ref{sec:warmup} the cases close by evaluation, and
\S\ref{sec:pushing} explains what replaces evaluation for the
symplectic group.

\subsection{Presentation construction theorems}\label{sec:constructions}

Relation sets over a disjoint union of alphabets are built from one
primitive, a three-way join $\Gamma_1 \aid{⋄} \Gamma_2 \aid{⋄}
\Gamma_3$ whose \emph{mixed} component $\Gamma_3 : \aid{WRel}\;(A
\uplus B)$ is the design parameter.  For the semidirect product the
mixed relation $\aid{ConjRelʷ}\;\mathit{conj}$ says that moving a
letter $h$ of the acting group past a letter $n$ of the normal subgroup
conjugates $n$, by a \emph{word}: $h \bullet n = \mathit{conj}\;h\;n
\bullet h$.  The semidirect presentation theorem takes the two relation
sets and the action, two well-definedness conditions on the action ---
it respects the acting group's axioms in its first argument and the
normal subgroup's in its second --- and presentations of the factors,
and returns a presentation of the semidirect product of the two
\emph{semantic} groups:

\begingroup\renewcommand{\AgdaCodeStyle}{\footnotesize}
\begin{code}
module _ (conj-hyph : ∀ {c d} n → c ===₂ d → (conj ʰ') c n ≈₁ (conj ʰ') d n)
         (conj-hypn : ∀ c {w v} → w ===₁ v → (conj ⁿ') c w ≈₁ (conj ⁿ') c v) where
  module Presentation (G1 G2 : Group 0ℓ 0ℓ)
    (p1 : Γ IsPresentationOf G1) (p2 : Δ IsPresentationOf G2) where
    dpres : (Γ ⋄ Δ ⋄ ConjRelʷ conj) IsPresentationOf G1⋊G2
\end{code}
\endgroup
\begin{code}[hide]
    dpres = proof-elided

extension-presentation : Set → Set → Set → Set → Set
\end{code}

\noindent Here \aid{G1⋊G2} is built by the library's standard-library
style semidirect product from an \aid{Action} record whose action is the
syntactic \aid{conj} transported through the two presentation
isomorphisms; its normal form is the product of the factors' normal
forms, and uniqueness is lifted from the factors by
\aid{by-completeness}.  A second mixed component, \aid{RelTwist},
twists each relator $u = v$ of a quotient by a correction word in the
normal subgroup, and
\begingroup\renewcommand{\AgdaCodeStyle}{\footnotesize}
\begin{code}
extension-presentation S R conj corr = S ⋄ EmptyRel ⋄ (ConjRelʷ conj ∪ RelTwist R corr)
\end{code}
\endgroup
is the recipe of Proposition~2.55 of the Handbook of Computational Group
Theory~\cite{holt2005handbook}, which the qupit paper invokes for its
Section~4.2: the relations of the normal subgroup, conjugation rules,
and each relator of the quotient corrected by the element by which its
lift fails to hold.  The library proves it as a presentation theorem for
any group extension realised by a short exact sequence and, in the form
used in \S\ref{sec:scalars}, for central extensions built from a
cocycle.  The semidirect product is the special case with trivial
corrections.  Finally, $\Zp$ itself is $\mathsf{Fin}\;p$ with a ring
solver, B\'ezout inversion, and Fermat's little theorem; a primitive
root $g$ and the witness that every unit is a power of it are
parameters, from which the discrete logarithm and the exponent law
$g^{k} = g^{k \bmod (p-1)}$ are derived.
